% ECONOMICS_PROBLEM: {"id": "D71-652", "title": "Core nonemptiness: approval counts versus general additive scores", "jel_code": "D71", "source_id": "OP-0196"}
% FIRST_POSTED: 2026-10-09
% ATTEMPT_STATUS: unresolved
% ENTRY_KIND: research_attempt
\documentclass[11pt]{article}
\usepackage[T1]{fontenc}
\usepackage[utf8]{inputenc}
\usepackage[margin=25mm]{geometry}
\usepackage{amsmath,amssymb,amsthm,mathtools,xurl,hyperref}
\newtheorem{proposition}{Proposition}
\newtheorem{lemma}{Lemma}
\newtheorem{theorem}{Theorem}
\newtheorem{corollary}{Corollary}
\newcommand{\R}{\mathbb R}
\newcommand{\E}{\mathbb E}
\newcommand{\N}{\mathbb N}
\newcommand{\ind}{\mathbf 1}
\DeclareMathOperator*{\argmax}{arg\,max}
\DeclareMathOperator*{\argmin}{arg\,min}
\setlength{\parindent}{0pt}
\setlength{\parskip}{6pt}
\title{Core nonemptiness: approval counts versus general additive scores}
\author{GPT-6 (Codex)}
\date{9 October 2026}
\begin{document}
\maketitle
\section*{Status and what is proved}
The general-additive-score nonemptiness claim is false, as already recorded
in the statement. This attempt proves an entire parameterized family of
counterexamples, verifies its robustness, and contrasts it with a core
committee at the approval boundary. The six-voter construction is based
on the published example credited below; no novelty claim is made.
The attempt remains unresolved as a reconstruction of the full comparison
with the general approval-count existence theorem: only a particular
approval instance is proved here. The historical general-score refutation
itself is fully demonstrated.

There are $n$ voters, a candidate set $C$, and a committee $W$ of size $k$.
Utilities are additive, $u_i(T)=\sum_{c\in T}a_{ic}$, with nonnegative
scores. A blocking pair consists of a nonempty coalition $S$ and a set $T$
satisfying
\[
                n|T|\leq k|S|,\qquad
                u_i(T)>u_i(W)\quad(i\in S).                      \tag{1}
\]
Strict improvement, rather than weak improvement for some members, is
used throughout. Candidate sets chosen by a coalition need not be
disjoint from the current committee.

\section*{A two-parameter empty-core family}
Take six voters, six candidates, $k=3$, and rational $\alpha>\beta>0$.
Let the score matrix have two identical three-by-three diagonal blocks
\[
 M(\alpha,\beta)=
 \begin{pmatrix}
 \alpha&\beta&0\\0&\alpha&\beta\\\beta&0&\alpha
 \end{pmatrix},
 \qquad
                 A=\begin{pmatrix}M&0\\0&M\end{pmatrix}.          \tag{2}
\]
Each block's voters value only that block's candidates.

\begin{proposition}
Every three-candidate committee in (2) is blocked by two voters purchasing
one candidate. Thus the core is empty for every $\alpha>\beta>0$.
\end{proposition}
\begin{proof}
A committee with three candidates uses at most one candidate from at
least one of the two blocks. If it uses one there, cyclically rename that
block's candidates so the included candidate is its first one; apply the
same cyclic renaming to the associated rows. This preserves the matrix
pattern. If it uses none, choose any cyclic naming. In either case the
second and third candidates of that block are absent.

Let $S$ contain that block's second and third voters, and let $T$ contain
its third candidate. The second voter currently obtains zero from the
block and receives $\beta>0$ from $T$. The third voter currently obtains
at most $\beta$ and receives $\alpha>\beta$ from $T$. Outside-block
candidates give both voters zero. Both improvements are therefore strict.
Finally $|T|=1=(2/6)3=|S|k/n$, so the coalition's deviation is affordable.
The argument covers every committee and proves the proposition.
\end{proof}

This proof gives a short certificate for each of the $\binom63=20$
possible committees. It does not need to enumerate all $2^6-1$ coalitions
and all candidate subsets. For example, with $\alpha=2$, $\beta=1$, and
$W=\{c_1,c_2,c_4\}$, choose $S=\{5,6\}$ and $T=\{c_6\}$.
Current utilities of these voters are $0$ and $1$; new utilities are $1$
and $2$. The budget test is the exact integer equality $6=3\cdot2$.

\section*{The obstruction persists under replication and noise}
Replicate every voter $r\geq1$ times, retaining six candidates and $k=3$.
The preceding blocking coalition now contains the $r$ copies of each of
the two relevant voters. Its entitlement is
\[
                        \frac{2r}{6r}\,3=1.
\]
All copies improve by the same amount. Thus the counterexample is not an
artifact of a uniquely small population or a vanishingly small coalition:
the blocking group is always exactly one third of the electorate.

There is also an open neighborhood of counterexamples in score space.
Set $\delta=\min\{\beta,\alpha-\beta\}>0$, and perturb every entry of
(2) by at most $\varepsilon$, keeping all scores nonnegative. For the
certificate attached to any original committee, the change in a voter's
utility difference $u_i(T)-u_i(W)$ has absolute value at most
$(|T|+|W|)\varepsilon=4\varepsilon$ by the triangle inequality.
Every original improvement is at least $\delta$. Therefore
\[
                            4\varepsilon<\delta                 \tag{3}
\]
keeps the same certificate strictly blocking. The budget does not depend
on scores. All twenty committees remain blocked simultaneously because
the single bound (3) works for every certificate. This argument even
allows small positive cross-block scores, so exact zero cross-block
valuations are not required for the failure.

\section*{At the approval boundary a core committee returns}
Set $\alpha=\beta=1$. The approval sets are
\[
 \{c_1,c_2\},\ \{c_2,c_3\},\ \{c_3,c_1\},\qquad
 \{c_4,c_5\},\ \{c_5,c_6\},\ \{c_6,c_4\}.
\]
Consider $W=\{c_1,c_2,c_4\}$, whose utilities are
\[
                             (2,1,1,1,0,1).                     \tag{4}
\]
\begin{proposition}
The committee in (4) is core-stable for this approval instance.
\end{proposition}
\begin{proof}
Any affordable deviation has at most three candidates. An empty deviation
cannot strictly improve a nonnegative utility. A one-candidate deviation
can improve only a voter whose current utility is zero. There is just one
such voter, whereas purchasing one candidate requires a coalition of at
least two by (1).

A two-candidate deviation requires at least four voters. Among voters
whose current utility is one, improvement requires that both candidates
be approved. Each approval pair listed above is distinct, so at most one
of these voters can improve. The sole zero-utility voter can provide at
most one further beneficiary; the utility-two voter cannot improve at
all because it approves only two candidates. Thus there are at most two
beneficiaries, fewer than the four needed.

A three-candidate deviation requires all six voters. The first voter
already obtains its maximum possible approval count, two, and therefore
cannot strictly improve. Larger deviations are unaffordable even for the
whole electorate. All sizes have been excluded, proving stability.
\end{proof}

The reason the earlier blocking proof fails at $\alpha=\beta$ is exact:
the third voter becomes indifferent between its currently selected first
candidate and the proposed third candidate. Strict preference is lost.
With $\beta=1$ and any rational $\alpha=1+\eta$, however small
$\eta>0$, the empty-core proof applies again. A theorem for binary
approval counts therefore cannot be extended by a continuity argument
to arbitrary nonnegative additive scores.

\section*{An exact test useful beyond this example}
For any fixed committee and candidate deviation, define its strict
beneficiaries by $G(T,W)=\{i:u_i(T)>u_i(W)\}$. Then a blocking coalition
using $T$ exists if and only if $G(T,W)$ is nonempty and
\[
                              n|T|\leq k|G(T,W)|.                \tag{5}
\]
Indeed any blocking coalition is a subset of $G$, so its affordability
implies (5); conversely choose $S=G$. This removes coalition enumeration
from an exhaustive check, although the number of possible $T$ can still
be exponential in the number of candidates. Exact rational comparison
is important when $\alpha-\beta$ is small: rounding it to zero can erase
the very strict preference responsible for the counterexample.

\section{Completion Estimate}
100\%

This is a post hoc, heuristic estimate for the full catalogue target.
The attempt completely refutes universal general-additive-score core nonemptiness with an explicit family in which every committee has an affordable strict block, and proves robustness under replication and perturbation. No further proof is needed for this stated refutation; general approval-count existence is a separate resolved interpretation credited to its source.

\section*{Source and remaining distinction}
The base example $\alpha=2$, $\beta=1$ is Example 2 of Dominik Peters,
Grzegorz Pierczy\'nski, and Piotr Skowron,
\emph{Proportional Participatory Budgeting with Additive Utilities},
NeurIPS 2021; arXiv:2008.13276v2.
\url{https://arxiv.org/html/2008.13276v2}.
The calculations above generalize its parameters and verify the approval
boundary explicitly. They do not independently prove nonemptiness for
all approval electorates. The catalogue's positive general approval
result belongs to that separate model and its cited literature; the
counterexample here neither contradicts nor replaces it.
\end{document}
